\documentclass[aspectratio=169]{beamer}

% ── Encodage / langue ────────────────────────────────────────
\usepackage[T1]{fontenc}
\usepackage[utf8]{inputenc}
\usepackage[french]{babel}
\usepackage{lmodern}
\usepackage{microtype}
\usepackage[backend=biber, style=authoryear]{biblatex}
\addbibresource{biblio.bib}
% ── Maths / figures ──────────────────────────────────────────
\usepackage{amsmath,amssymb,mathtools}
\usepackage{graphicx}
\usepackage{xcolor}
\usepackage{tikz}
\usetikzlibrary{arrows.meta}

% ============================================================
% PALETTE SOBRE
% ============================================================
\definecolor{Ink}{RGB}{18,24,35}
\definecolor{Navy}{RGB}{9,39,84}
\definecolor{Steel}{RGB}{32,88,150}
\definecolor{Slate}{RGB}{82,96,115}
\definecolor{Line}{RGB}{204,211,222}
\definecolor{Fog}{RGB}{244,247,251}
\definecolor{Pearl}{RGB}{253,253,255}
\definecolor{Gold}{RGB}{145,105,28}
\definecolor{Crimson}{RGB}{155,35,45}
\definecolor{Teal}{RGB}{20,120,105}
\definecolor{Violet}{RGB}{100,55,155}   % ← ajout : diffusion brownienne
% Alias pour le bilan
\colorlet{Good}{Teal}
\colorlet{Bad}{Crimson}
\usepackage{eurosym}
% ============================================================
% THÈME BEAMER
% ============================================================
\usetheme{default}
\usefonttheme{serif}
\setbeamertemplate{navigation symbols}{}
\setbeamertemplate{headline}{}

\setbeamercolor{background canvas}{bg=Pearl}
\setbeamercolor{normal text}{fg=Ink,bg=Pearl}
\setbeamercolor{frametitle}{fg=Navy,bg=Pearl}

\setbeamerfont{frametitle}{size=\large,series=\bfseries,family=\rmfamily}
\setbeamerfont{title}{size=\Large,series=\bfseries,family=\rmfamily}
\setbeamerfont{footnote}{size=\tiny}
\setbeamercolor{title}{fg=Navy}
\setbeamertemplate{itemize item}{\small\textcolor{Steel}{\textbullet}}
\setbeamertemplate{itemize subitem}{\scriptsize\textcolor{Slate}{$\circ$}}
\setbeamertemplate{enumerate item}{\textcolor{Navy}{\bfseries\insertenumlabel.}}

\setlength{\parskip}{0.15em}

% ── Ligne sous le titre de frame ─────────────────────────────
\addtobeamertemplate{frametitle}{}{%
  \vspace{-0.35em}\textcolor{Line}{\rule{\linewidth}{0.5pt}}\vspace{-0.05em}}

% ── Pied de page minimal ─────────────────────────────────────
\setbeamertemplate{footline}{%
  \begin{beamercolorbox}[wd=\paperwidth,ht=0.42cm,dp=0.18cm]{footline}
    \hfill\textcolor{Slate}{\scriptsize\insertframenumber/\inserttotalframenumber}\hspace{0.45cm}
  \end{beamercolorbox}}
\setbeamercolor{footline}{bg=Pearl}

% ── Diapo de section ─────────────────────────────────────────
\AtBeginSection[]{%
  \begin{frame}[plain]
  \begin{tikzpicture}[remember picture,overlay]
    \fill[Fog] (current page.north west) rectangle (current page.south east);
    \fill[Navy] (current page.north west) rectangle ++(0.42cm,-\paperheight);
    \fill[Gold]  ([xshift=0.42cm,yshift=-0.55cm]current page.north west)
                 rectangle ++(0.12cm,-3.15cm);
  \end{tikzpicture}
  \hspace{0.72cm}\begin{minipage}{0.82\textwidth}
    \vspace{2.2cm}
    {\color{Slate}\normalsize Partie~\insertsectionnumber}\par
    \vspace{0.25cm}
    {\color{Navy}\LARGE\bfseries\insertsectionhead}\par
    \vspace{0.35cm}
    {\color{Line}\rule{0.55\linewidth}{0.55pt}}\par
    \vspace{0.35cm}
    {\color{Slate}\small\tableofcontents[currentsection,hideallsubsections]}
  \end{minipage}
  \end{frame}
}

% ============================================================
% COMMANDES MATHÉMATIQUES
% ============================================================
\newcommand{\R}{\mathbb{R}}
\newcommand{\E}{\mathbb{E}}
\newcommand{\PP}{\mathbb{P}}
\newcommand{\1}{\mathbf{1}}
\newcommand{\dd}{\,\mathrm{d}}
\newcommand{\norm}[1]{\left\|#1\right\|}
\newcommand{\rsp}{\rho_{\mathrm{sp}}}
\DeclareMathOperator{\Tr}{Tr}

% ── Notations réseau ─────────────────────────────────────────
\newcommand{\Cset}{\mathcal{C}}
\newcommand{\Pset}{\mathcal{P}}
\newcommand{\nC}{N_c}
\newcommand{\nP}{N_p}
\newcommand{\rhobarC}{\bar{\rho}^{C}}
\newcommand{\rhobarP}{\bar{\rho}^{P}}

% ── Mise en forme typographique ───────────────────────────────
\newcommand{\key}[1]{\textcolor{Navy}{\textbf{#1}}}
\newcommand{\kw}[1]{\textcolor{Steel}{\textbf{#1}}}
\newcommand{\gold}[1]{\textcolor{Gold}{\textbf{#1}}}
\newcommand{\risk}[1]{\textcolor{Crimson}{\textbf{#1}}}
\newcommand{\muted}[1]{{\small\color{Slate}#1}}
\newcommand{\pro}[1]{\textcolor{Teal}{\textbf{#1}}}
\newcommand{\con}[1]{\textcolor{Crimson}{\textbf{#1}}}

% ── Filets, encadrés ─────────────────────────────────────────
\newcommand{\sep}{%
  \vspace{0.15em}\par\textcolor{Line}{\rule{\linewidth}{0.4pt}}\vspace{0.15em}}

\newcommand{\takeaway}[1]{%
  \vspace{0.25em}\par\noindent
  \begin{tikzpicture}
    \node[inner sep=0pt,outer sep=0pt,text width=0.92\linewidth,align=left]
      (txt) {\small\color{Slate}#1};
    \draw[Gold,line width=2.2pt]
      ([xshift=-0.55em]txt.north west)--([xshift=-0.55em]txt.south west);
  \end{tikzpicture}}

\newcommand{\prosbox}[1]{%
  \vspace{0.2em}\par\noindent
  \begin{tikzpicture}
    \node[inner sep=0pt,outer sep=0pt,text width=0.92\linewidth,align=left]
      (txt) {\small\color{Slate}#1};
    \draw[Teal,line width=2.2pt]
      ([xshift=-0.55em]txt.north west)--([xshift=-0.55em]txt.south west);
  \end{tikzpicture}}

\newcommand{\consbox}[1]{%
  \vspace{0.2em}\par\noindent
  \begin{tikzpicture}
    \node[inner sep=0pt,outer sep=0pt,text width=0.92\linewidth,align=left]
      (txt) {\small\color{Slate}#1};
    \draw[Crimson,line width=2.2pt]
      ([xshift=-0.55em]txt.north west)--([xshift=-0.55em]txt.south west);
  \end{tikzpicture}}

% ── Image sécurisée ──────────────────────────────────────────
\newcommand{\safeimage}[3]{%
  \IfFileExists{#1}{\includegraphics[width=#2,height=#3,keepaspectratio]{#1}}{%
    \begin{tikzpicture}
      \draw[Line,thick,rounded corners=2pt,fill=Fog] (0,0) rectangle (11.2,5.3);
      \node[align=center,text=Slate] at (5.6,2.65)
        {Image manquante : \texttt{\detokenize{#1}}\\
         Placer le fichier dans le même dossier que le \texttt{.tex}};
    \end{tikzpicture}}}

\title{\Large Apprentissage non-supervisé pour l'estimation de champs de
déformation à la surface de matériaux soumis à des
contraintes mécaniques}
\author{Projet Recherche - Oscar Egreteau}
\date{Sous la supervision de F.Sur, LORIA}
% % ── Logos en haut à droite ───────────────────────────────────
% \newcommand{\logoshautdroit}{%
%   \begin{tikzpicture}[remember picture,overlay]
%     \node[anchor=north east, xshift=-0.35cm, yshift=-0.18cm]
%       at (current page.north east) {%
%         \includegraphics[height=0.5cm]{Mines(1).png}\hspace{0.15cm}%
%         \includegraphics[height=0.5cm]{IECL-logo(1).jpg}%
%       };
%   \end{tikzpicture}%
% }
% \addtobeamertemplate{frametitle}{}{\logoshautdroit}

\usepackage[utf8]{inputenx} % Allows UTF-8 input
\usepackage[T1]{fontenc}    % Uses T1 font encoding (has proper Euro mapping)
\usepackage{textcomp}      % Provides the \texteuro command
\usepackage{marvosym}
\begin{document}

\begin{frame}
    \titlepage
\end{frame}

\begin{frame}
    \frametitle{Sommaire}
    \tableofcontents
\end{frame}

\section{Introduction}
\begin{frame}
    \frametitle{Introduction}
    \begin{itemize}
    \item Analyse du flot optique pour estimer le champ de déformation de matériaux
    \item Apprentissage non supervisé car difficulté d'avoir une réalité terrain : pas possible de connaître le champ de déformation sur un matériau sans l'utilisation d'une méthode d'estimation numérique.
\end{itemize}
% \begin{figure}
%         \centering 
%         \includegraphics[height=4cm]{Images/intro.png}
%         \caption*{Un exemple de champ de déformation}
%     \end{figure}
\end{frame}

\section{Flot optique}
\begin{frame}
\frametitle{Flot optique}
    Le flot optique est un champ de vecteurs qui représente le déplacement des pixels entre deux images consécutives. Son calcul repose sur une hypothèse fondamentale, celle de la luminosité constante.\\
    \begin{figure}[h]
      \centering
        \includegraphics[height=4cm]{Images/govas.png}
        \caption{Deux images consécutives}
    \end{figure}
\end{frame}

\begin{frame}
\frametitle{Flot Optique}
    \begin{figure}[h]
      \centering
        \includegraphics[height=4cm]{Images/output.png}
        \caption{Flot optique correspondant}
    \end{figure}
\end{frame}

\begin{frame}
    \frametitle{Flot Optique}
    \framesubtitle{Définition}
    Le flot optique $\bold{u}=(u_1,u_2)$ est régi par l'équation suivante, obtenue en en dérivant l'hypothèse de la luminosité constante ($\forall t\geq 0, \ I(t,x_1(t),x_2(t))=I_0$) :
    \begin{align*}
        \dfrac{\partial I}{\partial t}+\nabla I\cdot \bold{u} &= \dfrac{\partial I}{\partial t}+\dfrac{\partial I}{\partial x}u_1+\dfrac{\partial I}{\partial t}u_2\\
        &=0
    \end{align*}
    Qui est un problème mal posé car les solutions $\bold{u}=(u_1,u_2)$ sont posées sur une droite : il en existe une infinité.
\end{frame}

\section{Estimation de flot optique par réseaux de neurones}
\begin{frame}
\frametitle{Estimation de flot optique par réseaux de neurones}
\framesubtitle{Introduction}
    \begin{figure}
        \includegraphics[width=0.7\textwidth]{Images/Réseaux flot/Général/encodeur.png}
        \vspace{0.5em}
        \includegraphics[width=0.7\textwidth]{Images/Réseaux flot/Général/decodeur.png}
    \end{figure}
\end{frame}

\begin{frame}
    \frametitle{Estimation de flot optique par réseaux de neurones}
    \framesubtitle{Dataset}
\begin{figure}[h]
    \centering
    \begin{minipage}{0.45\textwidth}
        \centering
        \includegraphics[height=4cm]{Images/Flying/00001_img1.jpg}
    \end{minipage}
    \hfill
    \begin{minipage}{0.45\textwidth}
        \centering
        \includegraphics[height=4cm]{Images/Flying/00001_img2.jpg}
    \end{minipage}
    \caption{Un exemple de FlyingChairs}
\end{figure}
Le dataset contient $22872$ images différentes.
\end{frame}

\begin{frame}
    \frametitle{Estimation de flot optique par réseaux de neurones}
    \framesubtitle{Dataset}
\begin{figure}[h]
    \centering
    \begin{minipage}{0.45\textwidth}
        \centering
        \includegraphics[height=4cm]{Images/Flying/flot.png}
        \caption{Flot optique correspondant}
    \end{minipage}
    \hfill
    \begin{minipage}{0.45\textwidth}
        \centering
        \includegraphics[height=4cm]{Images/Flying/color-scheme.png}
        \caption{Roue de couleur pour lire le flot}
    \end{minipage}
\end{figure}
\end{frame}

\subsection{Apprentissage supervisé}
\begin{frame}
\frametitle{Apprentissage supervisé}
\framesubtitle{Résultats}
\begin{figure}[h]
    \centering
    \includegraphics[height=7cm]{Images/Réseaux flot/Supervisé/total_loss.png}
\end{figure}
\end{frame}

\begin{frame}
\frametitle{Apprentissage supervisé}
\framesubtitle{Résultats}
\begin{figure}[h]
    \centering
    \begin{minipage}{0.45\textwidth}
        \centering
        \includegraphics[height=4cm]{Images/Réseaux flot/Supervisé/output.png}
        \caption{Résultat obtenu}
    \end{minipage}
    \hfill
    \begin{minipage}{0.45\textwidth}
        \centering
        \includegraphics[height=4cm]{Images/Réseaux flot/Général/groundtruth.png}
        \caption{Groundtruth correpondante}
    \end{minipage}
\end{figure}
Ces résultats correspondent à ceux du papier
\end{frame}

\begin{frame}
\frametitle{Estimation de flot optique par réseaux de neurones}
\framesubtitle{Résultats}
Afin d'avoir une métrique de mesure pour qualifier les bons résultats ou non de notre réseau on utilise l'EPE:
\[ \sum\|V_i^{pred}-V_i^{gt}\|\]
\begin{itemize}
  \item EPE pour la méthode supervisée : $\text{EPE} = 3,89\text{ pixels}$
\end{itemize}
\end{frame}

\begin{frame}
\frametitle{Apprentissage supervisé}
\framesubtitle{Comparaison}
\begin{figure}[h]
    \centering
    \begin{minipage}{0.45\textwidth}
        \centering
        \includegraphics[height=4cm]{Images/Réseaux flot/Supervisé/output.png}
        \caption{Résultat obtenu}
    \end{minipage}
    \hfill
    \begin{minipage}{0.45\textwidth}
        \centering
        \includegraphics[height=4cm]{Images/comptvl.png}
        \caption{Résultat obtenu avec la méthode TV-L$1$}
    \end{minipage}
\end{figure}
\end{frame}

\begin{frame}
\frametitle{Apprentissage supervisé}
\framesubtitle{Comparaison}
On obtient de plus une différence notable entre les deux, le temps de calcul:
\begin{itemize}
  \item Inférence sur le set : instantané ($0.3$s)
  \pause
  \item Méthode TV-L$1$ : plutôt de l'ordre de la seconde ($1.2$s)
\end{itemize}
Une fois le set entrainé : plus efficace de faire l'inférence.\\
Mais il faut prendre en compte le temps d'entraînement. 
\end{frame}



\subsection{Apprentissage non-supervisé}
\begin{frame}
\frametitle{Apprentissage non-supervisé}
\framesubtitle{Principe}
    On remplace la loss $AEE$, par une nouvelle loss : 
    \[\mathcal L(u,v,I(x,y,t),I(x,y,t+1))=\ell_{\text{photometric}}(u,v,I(x,y,t),I(x,y,t+1))\]
    \[+\lambda \ell_{\text{smoothness}}(u,v)\] 
    Où : \\
    \begin{itemize}
        \item $u,v\in\mathbb R^{H\times W}$  sont les composantes horizontales et verticales du flot prédit.
        \item $\ell_{\text{photometric}}(u,v,I(x,y,t),I(x,y,t+1))=\displaystyle\sum_{i,j}\rho_D(I(i,j,t)-I(i+u_{i,j},j+v_{i,j},t+1))$
        \item $\ell_{\text{smoothness}}(u,v)=\displaystyle\sum_{j}^{H}\sum_{i}^{W}(\rho_S(u_{i,j}-u_{i+1,j}) + \rho_S(u_{i,j}-u_{i,j+1})+\rho_S(v_{i,j}-v_{i+1,j}) + \rho_S(v_{i,j}-v_{i,j+1}))$
    \end{itemize}
    
\end{frame}

\begin{frame}
    \frametitle{Estimation de flot optique par réseaux de neurones}
    \framesubtitle{Dataset non-supervisé}
\begin{figure}[h]
    \centering
    \begin{minipage}{0.45\textwidth}
        \centering
        \includegraphics[height=4cm]{Images/Flying/00001_img1.jpg}
    \end{minipage}
    \hfill
    \begin{minipage}{0.45\textwidth}
        \centering
        \includegraphics[height=4cm]{Images/Flying/00001_img2.jpg}
    \end{minipage}
    \caption{Un exemple de FlyingChairs}
\end{figure}
On garde les deux images consécutives mais on retire le flot réel pour l'entraînement.
\end{frame}

\begin{frame}
\frametitle{Apprentissage non-supervisé}
\framesubtitle{Résultats}
\begin{figure}[h]
    \centering
    \includegraphics[height=7cm]{Images/Réseaux flot/Non Supervisé/loss unsup.png}
\end{figure}
\end{frame}

\begin{frame}
\frametitle{Apprentissage non-supervisé}
\framesubtitle{Résultats}
\begin{figure}[h]
    \centering
    \begin{minipage}{0.45\textwidth}
        \centering
        \includegraphics[height=4cm]{Images/Réseaux flot/Non Supervisé/inf.png}
        \caption{Résultat obtenu}
    \end{minipage}
    \hfill
    \begin{minipage}{0.45\textwidth}
        \centering
        \includegraphics[height=4cm]{Images/Réseaux flot/Général/groundtruth.png}
        \caption{Groundtruth correpondante}
    \end{minipage}
\end{figure}
\end{frame}

\begin{frame}
\frametitle{Estimation de flot optique par réseaux de neurones}
\framesubtitle{Résultats}
Afin d'avoir une métrique de mesure pour qualifier les bons résultats ou non de notre réseau on utilise l'EPE:
\[ \sum\|V_i^{pred}-V_i^{gt}\|\]
\begin{itemize}
  \item EPE pour la méthode non-supervisée : $\text{EPE} = 16,65\text{ pixels}$
\end{itemize}
\end{frame}

\begin{frame}
\frametitle{Apprentissage non-supervisé}
\framesubtitle{Recherche du problème}
A la recherche du problème, chacun de ces points ont été explorés :
\begin{itemize}
  \item Architecture du réseau
  \pause
  \item Pipeline d'importation des images
  \pause
  \item Warping
  \pause
  \item Loss
\end{itemize}

\end{frame}

\begin{frame}
\frametitle{Apprentissage non-supervisé}
\framesubtitle{De la difficulté de la reproductibilité des résultats}
Les résultats du papier sont largement meilleurs, mais de tels résultats sont difficiles à reproduire pour diverses raisons:
\begin{itemize}
  \item Difficulté de tester différents hyperparamètres de part la durée d'entraînement
  \pause
  \item Détails possiblement non mentionnés dans le papier
  \pause
  \item Existence d'un code source, mais avec des dépendances obsolètes
  \pause
  \item Aucune mention de la seed initiale
  \pause
  \item Coût élevé
  \begin{itemize}
    \item Des composants
    \pause
    \item De l'électricité
    \pause
    \begin{itemize}
      \item $0,25\cdot0,500\cdot160=20\EUR$
    \end{itemize}
  \end{itemize}
\end{itemize}


\end{frame}

\begin{frame}
\frametitle{Estimation de flot optique par réseaux de neurones}
\framesubtitle{Conclusion}
\begin{itemize}
        \item Meilleurs résultats du supervisé : $\text{EPE} = 3,89$ vs $16,65$ px  sur les grands déplacements du dataset FlyingChairs.
        \pause
        \item Limites du non-supervisé :
        \pause
        \begin{itemize}
            \item Sensibilité extrême aux minima locaux de la fonction de coût photométrique.
            \pause
            \item Difficulté de convergence sans initialisation ou régularisation particulières.
            \pause
        \end{itemize}
        \item Toutefois les déplacements réels en mécanique des milieux continus sont sub-pixel ou de faible amplitude, on peut espérer que l'apprentissage non-supervisé fonctionnerait.
    \end{itemize}
\end{frame}

\section{Estimation de champs de déformation par réseaux de neurones}
\begin{frame}
\frametitle{Estimation de champs de déformation par réseaux de neurones}
\framesubtitle{Introduction}

\begin{itemize}
        \item Spécificité de la mécanique des matériaux : Les images étudiées présentent des textures de surfaces mouchetées.
        \pause
        \item Nature des déplacements : Contrairement au flot optique général, les champs de déformation mécanique ($\varepsilon_{xx}, \varepsilon_{yy}, \varepsilon_{xy}$) :
        \pause
        \begin{itemize}
            \item Sont de très faible amplitude (souvent sub-pixel).
            \pause
            \item Présentent de fortes contraintes de continuité spatiale (lois de comportement).
        \end{itemize}
        \pause
        \item Le défi de la vérité terrain : pas d'accès à la groundtruth donc on génère les champs de déplacement par ordinateur de façon analogue à FlyingChairs.
    \end{itemize}
\end{frame}

\subsection{Apprentissage supervisé}
\begin{frame}
\frametitle{Apprentissage supervisé}
\framesubtitle{Principe}
\begin{itemize}
  \item Même architecture que le FlowNet
  \item Images synthétiques de surfaces mouchetées, avec des deformations connues
\end{itemize}
\begin{figure}[h]
  \centering
  \includegraphics[height=4cm]{Images/Réseaux déformations/Sup/Star.png}
  \caption{Exemple d'une image du dataset}
\end{figure}
\end{frame}

\begin{frame}
\frametitle{Apprentissage supervisé}
\framesubtitle{Résultats}
Réseau déjà entrainé disponible
\begin{figure}[h]
    \centering
    \begin{minipage}{0.45\textwidth}
        \centering
        \includegraphics[height=2cm]{Images/Réseaux déformations/Sup/StrainNet-l.png}
        \caption{Déplacement prédit par le réseau}
    \end{minipage}
    \hfill
    \begin{minipage}{0.45\textwidth}
        \centering
        \includegraphics[height=2cm]{Images/Réseaux déformations/Sup/Reference.png}
        \caption{La groundtruth correspondante}
    \end{minipage}
\end{figure}
\begin{itemize}
  \item Prédiction globalement bonne sauf dans la partie gauche $\longrightarrow$ déplacements à très haute fréquence.
\end{itemize}
\end{frame}

\begin{frame}
\frametitle{Apprentissage supervisé}
\framesubtitle{Comparaison avec le FlowNet supervisé}
\begin{figure}
  \centering
  \includegraphics[height=5cm]{Images/Réseaux déformations/Sup/flownetsup.png}
  \caption{Déplacement prédit par les poids du FlowNet}
\end{figure}
\end{frame}

\begin{frame}
\frametitle{Apprentissage supervisé}
\framesubtitle{Comparaison avec le FlowNet supervisé}
Pistes d'amélioration:
\begin{itemize}
  \item Réseau inadapté car a appris sur des formes géométriques
  \pause
  \item Finetuning des poids déjà existants
\end{itemize}
\end{frame}

\subsection{Apprentissage non-supervisé}
\begin{frame}
\frametitle{Apprentissage non-supervisé}
\framesubtitle{Principe}
On reprend le principe du FlowNet non supervisé: 
\begin{itemize}
  \item On garde l'architecture du réseau
  \item On remplace la loss par 
\end{itemize}
\[\mathcal L(u,v,I(x,y,t),I(x,y,t+1))=\ell_{\text{photometric}}(u,v,I(x,y,t),I(x,y,t+1))\]
\[+\lambda \ell_{\text{smoothness}}(u,v)\]
\newline
$\longrightarrow$ A une chance de fonctionner car les déplacements sont plus petits que sur le set FlyingChairs.
\end{frame}

\begin{frame}
  \frametitle{Apprentissage non-supervisé}
\framesubtitle{Résultats}
\begin{figure}
  \includegraphics[height=4.5cm]{Images/Réseaux déformations/Unsup/unsup.png}
  \caption{Champ de déplacement prédit par le StrainNet non-supervisé}
\end{figure}
On obtient de plus une loss constante au cours de l'apprentissage.
\end{frame}

\begin{frame}
\frametitle{Apprentissage non-supervisé}
Il existe certaines pistes d'explication des résultats : 
\begin{itemize}
  \item Modification des hyperparamètres (une seule configuration a pu être testée)
  \pause
  \item Report des problèmes du FlowNet non-supervisé
\end{itemize}
\framesubtitle{Idées d'amélioration}

\end{frame}

\begin{frame}
\frametitle{Estimation de champs de déformation par réseaux de neurones}
\framesubtitle{Conclusion}
\begin{itemize}
        \item Le réseau StrainNet donne d'excellents résultats globaux sur les surfaces mouchetées, validant l'apprentissage des caractéristiques géométriques du matériau.
        \pause
        \item Échec du non-supervisé : Le transfert direct de la loss photométrique du flot optique classique ne fonctionne pas de manière satisfaisante sur ce type de texture :
        \pause
        \begin{itemize}
            \item Le réseau stagne ou génère des champs de déplacement lissés aberrants.
            \pause
            \item Les micro-déplacements sub-pixel sont noyés dans le bruit de la loss.
        \end{itemize}
        % \item L'absence de régularisation basée sur la physique mécanique empêche le modèle auto-supervisé de converger vers une solution réaliste.
    \end{itemize}
\end{frame}

\section{Conclusion et retour sur les objectifs}
\subsection{Retour sur les objectifs}
\begin{frame}
    \frametitle{Retour sur les objectifs}
    Objectifs fixés au mi-parcours
    \begin{itemize}
        \item Implémentation en Python de la méthode d'apprentissage non supervisé du flot optique.
        \item Application de celle-ci aux matériaux
    \end{itemize}

\end{frame}

\subsection{Conclusion}
\begin{frame}
\frametitle{Conclusion}
\begin{itemize}
        \item Apport des réseaux : L'inférence par deep learning offre un gain de temps de calcul considérable par rapport aux méthodes d'optimisation classiques telles que Horn-Schunck par exemple.
        \item Le verrou de l'auto-supervision : La perte photométrique seule peine à converger sur les textures mouchetées à haute fréquence sans un pré-entraînement supervisé.
    \end{itemize}

    \sep

    \prosbox{\textbf{Retour sur l'expérience en recherche :}
    \begin{itemize}
      \item Opportunité stimulante
      \item Goût nouveau pour la recherche
    \end{itemize}}
\end{frame}
\begin{frame}[allowframebreaks]
  \frametitle{Bibliographie}
  \nocite{*}  % add this before \printbibliography
  \printbibliography
\end{frame}
\end{document}
